% !PRINTED 530
\weberpage{530}
Dies $\varDelta$ ist eine durch $\Omega$ völlig bestimmte positive oder negative, aber niemals verschwindende ganze rationale Zahl, und es giebt immer eine aus ganzen Zahlen $\omega_1, \omega_2, \ldots, \omega_n$ bestehende Basis von $\Omega$, deren Discriminante gleich $\varDelta$ ist.

Eine solche Basis wollen wir kurz eine \gesperrt{Minimalbasis von} $\Omega$ nennen.

Verstehen wir unter $k_1, k_2, \ldots, k_n$ irgend welche ganze rationale Zahlen, so ist jede Zahl von der Gestalt
\weq{1}{\omega = k_1\omega_1 + k_2\omega_2 + \cdots + k_n\omega_n}
eine ganze algebraische Zahl, und wir beweisen jetzt den fundamentalen Satz:

1. \gesperrt{Wenn $\omega_1, \omega_2, \ldots, \omega_n$ eine Minimalbasis ist, so sind in der Form (1) alle ganzen Zahlen des Körpers $\Omega$ enthalten.}

Da $\omega_1, \omega_2, \ldots, \omega_n$ eine Basis ist, so kann zunächst jede Zahl in $\Omega$ in der Form (1) dargestellt werden, wenn für $k_1, k_2, \ldots, k_n$ rationale Brüche zugelassen werden. Nehmen wir also an, es sei eine ganze Zahl $\omega$ in der Form
\weq{2}{\omega = \frac{k_1\omega_1 + k_2\omega_2 + \cdots + k_n\omega_n}{k}}
darstellbar, worin $k_1, k_2, \ldots, k_n, k$ ganze rationale Zahlen sind, so dass nicht alle $k_1, k_2, \ldots, k_n$ mit $k$ einen gemeinschaftlichen Theiler haben. Ist $p$ irgend eine in $k$ aufgehende natürliche Primzahl und $k = pk'$, so muss wenigstens einer der Coëfficienten $k_1, k_2, \ldots, k_n$ durch $p$ untheilbar sein. Es sei etwa $k_1$ durch $p$ nicht theilbar; dann lässt sich die ganze rationale Zahl $l$ so bestimmen, dass $l k_1 \equiv 1 \; (\mathrm{mod}\ p)$, oder $(l k_1 - 1)$ durch $p$ theilbar wird. Es folgt dann aus (2):
\weq{3}{l k' \omega - \frac{l k_1 - 1}{p}\omega_1 = \frac{\omega_1 + l k_2\omega_2 + \cdots + l k_n\omega_n}{p} = \omega_1',}
und $\omega_1'$ ist gleichfalls eine ganze algebraische Zahl. Setzen wir
\weq{4}{\omega_2' = \omega_2, \ldots, \omega_n' = \omega_n,}
so bilden die Zahlen $\omega_1', \omega_2', \ldots, \omega_n'$ eine \gesperrt{ganzzahlige Basis} von $\Omega$, weil sich die Zahlen $\omega_1, \omega_2, \ldots, \omega_n$ und folglich alle Zahlen $\omega$ linear durch $\omega_1', \omega_2', \ldots, \omega_n'$ ausdrücken lassen, und die Formel (11) des vorigen Paragraphen ergiebt
\weq{5}{\varDelta(\omega_1', \omega_2', \ldots, \omega_n') = \frac{1}{p^2}\varDelta(\omega_1, \omega_2, \ldots, \omega_n).}
